%% kapa: exam=1 round=37 year=2026 period=1 form=A subject=civil_law subject_ko=민법 seq=5 number=5
%% source: Q-Net 감정평가사 기출문제 (cst00309) artlSeq=5258924 file=raw/1008/5258924/2026년 제37회 감정평가사 1차 1교시.pdf
%% conversion: kordoc 4.24.51 -> tools/gen1.py (auto); review=proofread-sonnet; flags=none
\begin{question}{5}
민법상 법인의 기관에 관한 설명으로 옳은 것은? (다툼이 있으면 판례에 따름)
\begin{choices}
  \item 감사는 법인의 대표기관으로서 그 성명과 주소는 등기사항이다.
  \item 법인과 이사의 이익이 상반하는 경우, 임시이사를 선임하여야 한다.
  \item 이사는 주무관청의 허가가 없으면 사임할 수 없다.
  \item 정관에 이사의 해임사유에 관한 규정이 있는 경우, 특별한 사정이 없는 한 법인은 정관에서 정하지 아니한 사유로 이사를 해임할 수 없다.
  \item 법인은 대표권이 없는 이사의 불법행위에 대해서도 민법 제35조의 법인의 불법행위 책임을 진다.
\end{choices}
\end{question}
\answer{4}
